The complementary assignment regions and sharp assignment imply \(X_i\in\mathcal A_t\) if and only if \(T_i=t\).  The definitions of the empirical counts, local-constant fits, admissible sets, deterministic tie break, and fallback show that \(\widehat\tau_n\) is a single measurable map of \(\mathcal D_n\), the known geometry, \(L\), \(\sigma\), and \(n\).  In particular, it uses neither \(\kappa\) nor the score density.

Fix \(2<\kappa\le\overline\kappa\), let \(P\in\mathcal P_{\kappa}^{\mathrm{base}}\), and put
\[
 r=(\log n/n)^{1/(\kappa+2)}.
\]
For all sufficiently large \(n\), the dyadic grid in \(\mathrm{def:dyadic\text{-}grid}\) contains \(h_*\) satisfying \(r\le h_*<2r\).  Let \(z_1,\ldots,z_J\) be a maximal \(h_*/2\)-separated subset of \(\mathcal B\).  Maximality makes it an \(h_*/2\)-net, and the upper half of \(\mathrm{ass:global\text{-}boundary\text{-}entropy}\), inherited through \(\mathrm{def:base\text{-}thickness\text{-}class}\), gives
\[
 J\le 2C_{\mathcal B}h_*^{-1}. \tag{8}
\]
For every net center and side, \(\mathrm{ass:polynomial\text{-}lower\text{-}mass}\) gives
\[
 q_{t,P}(z_j,h_*/2)\ge c_m(h_*/2)^\kappa. \tag{9}
\]
Under \(\mathrm{ass:iid\text{-}sampling}\), the associated count is binomial.  If \(N\sim\operatorname{Bin}(n,q)\), Markov's inequality applied to \(2^{-N}\) yields
\[
 \Pr\{N\le nq/2\}
 \le 2^{nq/2}\mathbb E(2^{-N})
 =2^{nq/2}(1-q/2)^n
 \le e^{-nq/8}. \tag{10}
\]
Taking a union bound over (8), the two sides, and (9), and then using \(B(z_j,h_*/2)\subseteq B(x,h_*)\) whenever \(d(x,z_j)\le h_*/2\), produces a score-measurable event \(G_n\) on which
\[
 \inf_{x\in\mathcal B,\ t\in\{0,1\}}N_{t,n}(x,h_*)
 \ge c n c_mh_*^\kappa, \tag{11}
\]
and
\[
 \sup_{P\in\mathcal P_{\kappa}^{\mathrm{base}}}P(G_n^c)
 \le Ch_*^{-1}\exp\{-c n c_mh_*^\kappa\}. \tag{12}
\]
Because \(h_*\asymp r\),
\[
 \frac{nh_*^\kappa}{\log(e+n/h_*)}
 \ge c\left(\frac n{\log n}\right)^{2/(\overline\kappa+2)}\longrightarrow\infty \tag{13}
\]
uniformly over the declared exponent range.  Hence (11) makes \(h_*\) admissible for every \(x,t\) on \(G_n\), for all sufficiently large \(n\).  The minimizing property in \(\mathrm{def:count\text{-}adaptive\text{-}bandwidth}\), (11), and \(nr^{\kappa+2}=\log n\) imply
\[
 \begin{aligned}
 &L\widehat h_{t,n}(x)
 +\sigma\sqrt{\frac{\log(e+n/\widehat h_{t,n}(x))}
 {N_{t,n}(x,\widehat h_{t,n}(x))}}\\
 &\qquad\le Lh_*+
 \sigma\sqrt{\frac{\log(e+n/h_*)}{c n c_mh_*^\kappa}}
 \le Cr. \tag{14}
 \end{aligned}
\]

The lower-mass condition also supplies the support premise needed at the boundary.  Indeed, for any \(x\in\mathcal B\), side \(t\), and radius \(s>0\), taking \(b=\min\{h_0,s/2\}\) gives
\[
 P_X\{X\in\mathcal A_t:d(X,x)<s\}
 \ge q_{t,P}(x,b)\ge c_mb^\kappa>0. \tag{15}
\]
Thus \(x\in\mathcal X_{t,P}\).  The Lipschitz condition in \(\mathrm{ass:lipschitz\text{-}side\text{-}means}\) and the trace definition yield
\[
 \theta_{t,P}(x)=g_{t,P}(x). \tag{16}
\]
For a nonempty selected averaging set, every included score lies within \(\widehat h_{t,n}(x)\) of \(x\), so (16) bounds its conditional bias by \(L\widehat h_{t,n}(x)\).

Let \(Z_n\) be the maximum of the normalized residual sums over the nonempty disk patterns in \(\mathrm{lem:disk\text{-}pattern\text{-}sub\text{-}gaussian\text{-}control}\).  That lemma and \(|\mathcal H_n|\le 2+\log_2 n\) imply
\[
 \mathbb E_P(Z_n\mid X_1,\ldots,X_n)\le C\sqrt{\log n}. \tag{17}
\]
Because \(h\le h_0<1\) entails \(\log(e+n/h)\ge\log n\), equations (14), (16), and (17), followed by the triangle inequality over the two sides, give
\[
 \mathbb E_P\!\left[\sup_{x\in\mathcal B}|\widehat\tau_n(x)-\tau_P(x)|
 \ \middle|\ X_1,\ldots,X_n\right]\mathbf 1_{G_n}
 \le Cr\mathbf 1_{G_n}. \tag{18}
\]

On \(G_n^c\), a selected set is either nonempty or is the \(h_0\)-fallback.  If it is nonempty, (16), Lipschitzness, and \(\widehat h_{t,n}(x)\le h_0\) bound its conditional mean error by \(Lh_0\); if it is empty, the fit is zero and \(\mathrm{ass:bounded\text{-}side\text{-}traces}\) bounds the trace error by \(L\).  Since every nonempty count is at least one, (17) therefore gives the conditional bound \(C\{L+\sigma\sqrt{\log n}\}\) for the supremum loss.  The right-hand side of (12) decreases faster than any power uniformly in \(2<\kappa\le\overline\kappa\), by (13).  Integrating (18) and the off-event bound proves
\[
 \sup_{P\in\mathcal P_{\kappa}^{\mathrm{base}}}
 \mathbb E_P\sup_{x\in\mathcal B}|\widehat\tau_n(x)-\tau_P(x)|
 \le Cr. \tag{19}
\]

For the fallback probability, repeat the net argument with the fixed radius \(h_0/2\).  The net has at most \(2C_{\mathcal B}h_0^{-1}\) points, while every center-side count has expectation at least \(nc_m(h_0/2)^\kappa\), which uniformly exceeds \(16\log(e+n/h_0)\) for all sufficiently large \(n\).  Equation (10), a union bound, and ball inclusion imply
\[
 \sup_{P\in\mathcal P_{\kappa}^{\mathrm{base}}}
 P\{\exists x,t:\widehat{\mathcal H}_{t,n}(x)=\varnothing\}
 \le C\exp\{-c n c_mh_0^\kappa\}
 \le Cn^3\exp\{-c n c_mh_0^\kappa\}. \tag{20}
\]
The off-event calculation above already includes this event's contribution.

Finally, \(\mathcal P_{\kappa}^{\mathrm{perv}}\subseteq\mathcal P_{\kappa}^{\mathrm{base}}\) by \(\mathrm{def:pervasive\text{-}thinning\text{-}class}\).  The single estimator bound (19) gives both \(R_n(\mathcal P_{\kappa}^{\mathrm{base}})\le Cr\) and \(R_n(\mathcal P_{\kappa}^{\mathrm{perv}})\le Cr\).  The pervasive Gaussian packing converse in \(\mathrm{thm:thickness\text{-}entropy\text{-}separation}\), certified by \(\mathrm{lem:fano\text{-}tent\text{-}validity}\) and \(\mathrm{lem:finite\text{-}information\text{-}testing}\), gives \(R_n(\mathcal P_{\kappa}^{\mathrm{perv}})\ge cr\).  When the pervasive subclass is nonempty, monotonicity of minimax risk under class inclusion gives
\[
 R_n(\mathcal P_{\kappa}^{\mathrm{base}})
 \ge R_n(\mathcal P_{\kappa}^{\mathrm{perv}})\ge cr. \tag{21}
\]
Equations (19)--(21) establish both matched frontiers.